\IfFileExists{stacks-project.cls}{%
\documentclass{stacks-project}
}{%
\documentclass{amsart}
}

% For dealing with references we use the comment environment
\usepackage{verbatim}
\newenvironment{reference}{\comment}{\endcomment}
%\newenvironment{reference}{}{}
\newenvironment{slogan}{\comment}{\endcomment}
\newenvironment{history}{\comment}{\endcomment}

% For commutative diagrams we use Xy-pic
\usepackage[all]{xy}

% We use 2cell for 2-commutative diagrams.
\xyoption{2cell}
\UseAllTwocells

% We use multicol for the list of chapters between chapters
\usepackage{multicol}

% This is generally recommended for better output
\usepackage{lmodern}
\usepackage[T1]{fontenc}

% For cross-file-references
\usepackage{xr-hyper}

% Package for hypertext links:
\usepackage{hyperref}

% For any local file, say "hello.tex" you want to link to please
% use \externaldocument[hello-]{hello}
\externaldocument[introduction-]{introduction}
\externaldocument[conventions-]{conventions}
\externaldocument[sets-]{sets}
\externaldocument[categories-]{categories}
\externaldocument[topology-]{topology}
\externaldocument[sheaves-]{sheaves}
\externaldocument[sites-]{sites}
\externaldocument[stacks-]{stacks}
\externaldocument[fields-]{fields}
\externaldocument[algebra-]{algebra}
\externaldocument[brauer-]{brauer}
\externaldocument[homology-]{homology}
\externaldocument[derived-]{derived}
\externaldocument[simplicial-]{simplicial}
\externaldocument[more-algebra-]{more-algebra}
\externaldocument[smoothing-]{smoothing}
\externaldocument[modules-]{modules}
\externaldocument[sites-modules-]{sites-modules}
\externaldocument[injectives-]{injectives}
\externaldocument[cohomology-]{cohomology}
\externaldocument[sites-cohomology-]{sites-cohomology}
\externaldocument[dga-]{dga}
\externaldocument[dpa-]{dpa}
\externaldocument[sdga-]{sdga}
\externaldocument[hypercovering-]{hypercovering}
\externaldocument[schemes-]{schemes}
\externaldocument[constructions-]{constructions}
\externaldocument[properties-]{properties}
\externaldocument[morphisms-]{morphisms}
\externaldocument[coherent-]{coherent}
\externaldocument[divisors-]{divisors}
\externaldocument[limits-]{limits}
\externaldocument[varieties-]{varieties}
\externaldocument[topologies-]{topologies}
\externaldocument[descent-]{descent}
\externaldocument[perfect-]{perfect}
\externaldocument[more-morphisms-]{more-morphisms}
\externaldocument[flat-]{flat}
\externaldocument[groupoids-]{groupoids}
\externaldocument[more-groupoids-]{more-groupoids}
\externaldocument[etale-]{etale}
\externaldocument[chow-]{chow}
\externaldocument[intersection-]{intersection}
\externaldocument[pic-]{pic}
\externaldocument[weil-]{weil}
\externaldocument[adequate-]{adequate}
\externaldocument[dualizing-]{dualizing}
\externaldocument[duality-]{duality}
\externaldocument[discriminant-]{discriminant}
\externaldocument[derham-]{derham}
\externaldocument[local-cohomology-]{local-cohomology}
\externaldocument[algebraization-]{algebraization}
\externaldocument[curves-]{curves}
\externaldocument[resolve-]{resolve}
\externaldocument[models-]{models}
\externaldocument[functors-]{functors}
\externaldocument[equiv-]{equiv}
\externaldocument[pione-]{pione}
\externaldocument[etale-cohomology-]{etale-cohomology}
\externaldocument[proetale-]{proetale}
\externaldocument[relative-cycles-]{relative-cycles}
\externaldocument[more-etale-]{more-etale}
\externaldocument[trace-]{trace}
\externaldocument[crystalline-]{crystalline}
\externaldocument[spaces-]{spaces}
\externaldocument[spaces-properties-]{spaces-properties}
\externaldocument[spaces-morphisms-]{spaces-morphisms}
\externaldocument[decent-spaces-]{decent-spaces}
\externaldocument[spaces-cohomology-]{spaces-cohomology}
\externaldocument[spaces-limits-]{spaces-limits}
\externaldocument[spaces-divisors-]{spaces-divisors}
\externaldocument[spaces-over-fields-]{spaces-over-fields}
\externaldocument[spaces-topologies-]{spaces-topologies}
\externaldocument[spaces-descent-]{spaces-descent}
\externaldocument[spaces-perfect-]{spaces-perfect}
\externaldocument[spaces-more-morphisms-]{spaces-more-morphisms}
\externaldocument[spaces-flat-]{spaces-flat}
\externaldocument[spaces-groupoids-]{spaces-groupoids}
\externaldocument[spaces-more-groupoids-]{spaces-more-groupoids}
\externaldocument[bootstrap-]{bootstrap}
\externaldocument[spaces-pushouts-]{spaces-pushouts}
\externaldocument[spaces-chow-]{spaces-chow}
\externaldocument[groupoids-quotients-]{groupoids-quotients}
\externaldocument[spaces-more-cohomology-]{spaces-more-cohomology}
\externaldocument[spaces-simplicial-]{spaces-simplicial}
\externaldocument[spaces-duality-]{spaces-duality}
\externaldocument[formal-spaces-]{formal-spaces}
\externaldocument[restricted-]{restricted}
\externaldocument[spaces-resolve-]{spaces-resolve}
\externaldocument[formal-defos-]{formal-defos}
\externaldocument[defos-]{defos}
\externaldocument[cotangent-]{cotangent}
\externaldocument[examples-defos-]{examples-defos}
\externaldocument[algebraic-]{algebraic}
\externaldocument[examples-stacks-]{examples-stacks}
\externaldocument[stacks-sheaves-]{stacks-sheaves}
\externaldocument[criteria-]{criteria}
\externaldocument[artin-]{artin}
\externaldocument[quot-]{quot}
\externaldocument[stacks-properties-]{stacks-properties}
\externaldocument[stacks-morphisms-]{stacks-morphisms}
\externaldocument[stacks-limits-]{stacks-limits}
\externaldocument[stacks-cohomology-]{stacks-cohomology}
\externaldocument[stacks-perfect-]{stacks-perfect}
\externaldocument[stacks-introduction-]{stacks-introduction}
\externaldocument[stacks-more-morphisms-]{stacks-more-morphisms}
\externaldocument[stacks-geometry-]{stacks-geometry}
\externaldocument[moduli-]{moduli}
\externaldocument[moduli-curves-]{moduli-curves}
\externaldocument[examples-]{examples}
\externaldocument[exercises-]{exercises}
\externaldocument[guide-]{guide}
\externaldocument[desirables-]{desirables}
\externaldocument[coding-]{coding}
\externaldocument[obsolete-]{obsolete}
\externaldocument[fdl-]{fdl}
\externaldocument[index-]{index}

% Theorem environments.
%
\theoremstyle{plain}
\newtheorem{theorem}[subsection]{Theorem}
\newtheorem{proposition}[subsection]{Proposition}
\newtheorem{lemma}[subsection]{Lemma}

\theoremstyle{definition}
\newtheorem{definition}[subsection]{Definition}
\newtheorem{example}[subsection]{Example}
\newtheorem{exercise}[subsection]{Exercise}
\newtheorem{situation}[subsection]{Situation}

\theoremstyle{remark}
\newtheorem{remark}[subsection]{Remark}
\newtheorem{remarks}[subsection]{Remarks}

\numberwithin{equation}{subsection}

% Macros
%
\def\lim{\mathop{\mathrm{lim}}\nolimits}
\def\colim{\mathop{\mathrm{colim}}\nolimits}
\def\Spec{\mathop{\mathrm{Spec}}}
\def\Hom{\mathop{\mathrm{Hom}}\nolimits}
\def\Ext{\mathop{\mathrm{Ext}}\nolimits}
\def\SheafHom{\mathop{\mathcal{H}\!\mathit{om}}\nolimits}
\def\SheafExt{\mathop{\mathcal{E}\!\mathit{xt}}\nolimits}
\def\Sch{\mathit{Sch}}
\def\Mor{\mathop{\mathrm{Mor}}\nolimits}
\def\Ob{\mathop{\mathrm{Ob}}\nolimits}
\def\Sh{\mathop{\mathit{Sh}}\nolimits}
\def\NL{\mathop{N\!L}\nolimits}
\def\CH{\mathop{\mathrm{CH}}\nolimits}
\def\proetale{{pro\text{-}\acute{e}tale}}
\def\etale{{\acute{e}tale}}
\def\QCoh{\mathit{QCoh}}
\def\Ker{\mathop{\mathrm{Ker}}}
\def\Im{\mathop{\mathrm{Im}}}
\def\Coker{\mathop{\mathrm{Coker}}}
\def\Coim{\mathop{\mathrm{Coim}}}

% Boxtimes
%
\DeclareMathSymbol{\boxtimes}{\mathbin}{AMSa}{"02}

%
% Macros for moduli stacks/spaces
%
\def\QCohstack{\mathcal{QC}\!\mathit{oh}}
\def\Cohstack{\mathcal{C}\!\mathit{oh}}
\def\Spacesstack{\mathcal{S}\!\mathit{paces}}
\def\Quotfunctor{\mathrm{Quot}}
\def\Hilbfunctor{\mathrm{Hilb}}
\def\Curvesstack{\mathcal{C}\!\mathit{urves}}
\def\Polarizedstack{\mathcal{P}\!\mathit{olarized}}
\def\Complexesstack{\mathcal{C}\!\mathit{omplexes}}
% \Pic is the operator that assigns to X its picard group, usage \Pic(X)
% \Picardstack_{X/B} denotes the Picard stack of X over B
% \Picardfunctor_{X/B} denotes the Picard functor of X over B
\def\Pic{\mathop{\mathrm{Pic}}\nolimits}
\def\Picardstack{\mathcal{P}\!\mathit{ic}}
\def\Picardfunctor{\mathrm{Pic}}
\def\Deformationcategory{\mathcal{D}\!\mathit{ef}}
\usepackage{cleveref}
\crefname{originallemma}{original lemma}{original lemmas}
\Crefname{originallemma}{Original lemma}{Original lemmas}
\makeatletter
\newlabel{sites-cohomology-lemma-cech-vanish-collection}{{03F9}{source}{}{sourcepointer-1}{}}
\newlabel{sites-cohomology-lemma-cech-vanish-collection@cref}{{[originallemma][0][]03F9}{[1][0][]source}{}{}{}}
\makeatother
\crefname{originallemma}{original lemma}{original lemmas}
\Crefname{originallemma}{Original lemma}{Original lemmas}
\makeatletter
\newlabel{descent-lemma-standard-covering-Cech-quasi-coherent}{{03FJ}{source}{}{sourcepointer-2}{}}
\newlabel{descent-lemma-standard-covering-Cech-quasi-coherent@cref}{{[originallemma][0][]03FJ}{[1][0][]source}{}{}{}}
\makeatother
\crefname{originallemma}{original lemma}{original lemmas}
\Crefname{originallemma}{Original lemma}{Original lemmas}
\makeatletter
\newlabel{spaces-descent-proposition-fpqc-descent-quasi-coherent}{{04W8}{source}{}{sourcepointer-3}{}}
\newlabel{spaces-descent-proposition-fpqc-descent-quasi-coherent@cref}{{[originallemma][0][]04W8}{[1][0][]source}{}{}{}}
\makeatother
\crefname{originallemma}{original lemma}{original lemmas}
\Crefname{originallemma}{Original lemma}{Original lemmas}
\makeatletter
\newlabel{sites-lemma-mapping-property-glue}{{04TS}{source}{}{sourcepointer-4}{}}
\newlabel{sites-lemma-mapping-property-glue@cref}{{[originallemma][0][]04TS}{[1][0][]source}{}{}{}}
\makeatother
\crefname{originallemma}{original lemma}{original lemmas}
\Crefname{originallemma}{Original lemma}{Original lemmas}
\makeatletter
\newlabel{etale-cohomology-lemma-review-quasi-coherent}{{0DEW}{source}{}{sourcepointer-5}{}}
\newlabel{etale-cohomology-lemma-review-quasi-coherent@cref}{{[originallemma][0][]0DEW}{[1][0][]source}{}{}{}}
\makeatother
\crefname{originallemma}{original lemma}{original lemmas}
\Crefname{originallemma}{Original lemma}{Original lemmas}
\makeatletter
\newlabel{sites-modules-section-local}{{03DK}{source}{}{sourcepointer-6}{}}
\newlabel{sites-modules-section-local@cref}{{[originallemma][0][]03DK}{[1][0][]source}{}{}{}}
\makeatother
\crefname{originallemma}{original lemma}{original lemmas}
\Crefname{originallemma}{Original lemma}{Original lemmas}
\makeatletter
\newlabel{sites-definition-special-cocontinuous-functor}{{03CG}{source}{}{sourcepointer-7}{}}
\newlabel{sites-definition-special-cocontinuous-functor@cref}{{[originallemma][0][]03CG}{[1][0][]source}{}{}{}}
\makeatother
\crefname{originallemma}{original lemma}{original lemmas}
\Crefname{originallemma}{Original lemma}{Original lemmas}
\makeatletter
\newlabel{sites-lemma-equivalence}{{03A0}{source}{}{sourcepointer-8}{}}
\newlabel{sites-lemma-equivalence@cref}{{[originallemma][0][]03A0}{[1][0][]source}{}{}{}}
\makeatother
\title{Vanishing for the homology sheaf of a complex of quasi-coherent modules on the etale and fppf sites}
\author{The Stacks Project Authors}
\date{stacks-project:spaces-more-cohomology.tex:local-snapshot-20261012; GNU FDL 1.2 or later}
\begin{document}
\maketitle
\noindent\textit{Editorial note.} The counted statement is the author's own lemma environment labelled lemma-vanishing-adequate (source lines 1615-1639); the proof environment beginning after the author's following remark paragraph (source lines 1641-1647) is its own proof and closes at source line 1703.\par\medskip
\noindent\textit{Editorial note.} The remark between the statement and the proof (introducing adequate modules and pointing to the Adequate Modules chapter) is part of the surrounding narrative, not of the counted statement or proof, and is not retained.\par\medskip
\noindent\textit{Editorial note.} The proof invokes one result of the same chapter by label (retained as context with the generator's transitive expansion) and two cross-chapter results (Cohomology on Sites, Descent), each declared as a tagged source pointer.\par\medskip
\noindent\textit{Editorial note.} The retained preamble is the chapter's own (sources/stacks-spaces-more-cohomology/preamble.tex, taken with its real line count), a verbatim excerpt of the native Xy-pic setup and theorem declarations; no macro was invented or redefined.\par\medskip
\noindent\textit{Editorial scope.} One original result of the Stacks Project's chapter Cohomology of Algebraic Spaces, counted once: the vanishing of the etale and fppf cohomology of the homology sheaf of a complex of quasi-coherent modules at affine objects.\par\medskip
\noindent\textit{Allowed standard prerequisites (editorial).}\begin{itemize}
\item Restriction of the big etale site to a small one. For f : U -> X in the big etale site one has i\_f o pi\_X = f as morphisms of ringed sites U\_etale -> X\_etale, and restriction along i\_f is exact.. It identifies H\_etale|\_\{U\_etale\} with the degree-one homology of f\textasciicircum{}*F\_bullet at source lines 1650-1658.
\item The quasi-coherent review lemma. The presheaf Y -> Gamma(Y, f\textasciicircum{}*F) satisfies the fpqc (hence fppf and etale) sheaf condition (the author's lemma-review-quasi-coherent, retained as context).. It gives that H\_etale is an fppf sheaf and hence H\_etale = epsilon\_\{X,*\}H\_fppf at source lines 1672-1682.
\item Cech vanishing criterion for a collection of coverings. If Cech cohomology vanishes for a suitable collection of coverings, then the sheaf cohomology vanishes (Cohomology on Sites, declared as a tagged source pointer).. It reduces the vanishing to the displayed Cech comparison at source lines 1685-1691.
\item Cech vanishing for quasi-coherent modules on a standard covering. The standard covering Cech complex of a quasi-coherent module has no higher cohomology (Descent, declared as a tagged source pointer).. It finishes the proof at source lines 1700-1702.
\end{itemize}\medskip
\ifdefined\hypertarget\hypertarget{sourcepointer-1}{}\fi
\noindent\textit{Source reference.} Original lemma 03F9: sources/stacks-chow/tags.txt, source lines 4484--4484. Cross-chapter pointer whose target carries the canonical tag 03F9 in the project tag list; declared so the reference resolves inside the assembled package, while every step the counted argument uses remains inside the retained spans.\par\medskip
\ifdefined\hypertarget\hypertarget{sourcepointer-2}{}\fi
\noindent\textit{Source reference.} Original lemma 03FJ: sources/stacks-chow/tags.txt, source lines 4494--4494. Cross-chapter pointer whose target carries the canonical tag 03FJ in the project tag list; declared so the reference resolves inside the assembled package, while every step the counted argument uses remains inside the retained spans.\par\medskip
\ifdefined\hypertarget\hypertarget{sourcepointer-3}{}\fi
\noindent\textit{Source reference.} Original lemma 04W8: sources/stacks-chow/tags.txt, source lines 6351--6351. Cross-chapter pointer whose target carries the canonical tag 04W8 in the project tag list; declared so the reference resolves inside the assembled package, while every step the counted argument uses remains inside the retained spans.\par\medskip
\ifdefined\hypertarget\hypertarget{sourcepointer-4}{}\fi
\noindent\textit{Source reference.} Original lemma 04TS: sources/stacks-chow/tags.txt, source lines 6265--6265. Cross-chapter pointer whose target carries the canonical tag 04TS in the project tag list; declared so the reference resolves inside the assembled package, while every step the counted argument uses remains inside the retained spans.\par\medskip
\ifdefined\hypertarget\hypertarget{sourcepointer-5}{}\fi
\noindent\textit{Source reference.} Original lemma 0DEW: sources/stacks-chow/tags.txt, source lines 16816--16816. Cross-chapter pointer whose target carries the canonical tag 0DEW in the project tag list; declared so the reference resolves inside the assembled package, while every step the counted argument uses remains inside the retained spans.\par\medskip
\ifdefined\hypertarget\hypertarget{sourcepointer-6}{}\fi
\noindent\textit{Source reference.} Original lemma 03DK: sources/stacks-chow/tags.txt, source lines 4423--4423. Cross-chapter pointer whose target carries the canonical tag 03DK in the project tag list; declared so the reference resolves inside the assembled package, while every step the counted argument uses remains inside the retained spans.\par\medskip
\ifdefined\hypertarget\hypertarget{sourcepointer-7}{}\fi
\noindent\textit{Source reference.} Original lemma 03CG: sources/stacks-chow/tags.txt, source lines 4383--4383. Cross-chapter pointer whose target carries the canonical tag 03CG in the project tag list; declared so the reference resolves inside the assembled package, while every step the counted argument uses remains inside the retained spans.\par\medskip
\ifdefined\hypertarget\hypertarget{sourcepointer-8}{}\fi
\noindent\textit{Source reference.} Original lemma 03A0: sources/stacks-chow/tags.txt, source lines 4291--4291. Cross-chapter pointer whose target carries the canonical tag 03A0 in the project tag list; declared so the reference resolves inside the assembled package, while every step the counted argument uses remains inside the retained spans.\par\medskip
\begin{lemma}
\label{lemma-review-quasi-coherent}
Let $S$ be a scheme. Let $X$ be an algebraic space over $S$.
Let $\mathcal{F}$ be a quasi-coherent $\mathcal{O}_X$-module.
\begin{enumerate}
\item The rule
$$
\mathcal{F}^a : (\textit{Spaces}/X)_\etale \longrightarrow \textit{Ab},\quad
(f : Y \to X) \longmapsto \Gamma(Y, f^*\mathcal{F})
$$
satisfies the sheaf condition for fpqc and a fortiori
fppf and \'etale coverings,
\item $\mathcal{F}^a = \pi_X^*\mathcal{F}$ on $(\textit{Spaces}/X)_\etale$,
\item $\mathcal{F}^a = a_X^*\mathcal{F}$ on $(\textit{Spaces}/X)_{fppf}$,
\item the rule $\mathcal{F} \mapsto \mathcal{F}^a$ defines
an equivalence between quasi-coherent $\mathcal{O}_X$-modules
and quasi-coherent modules on
$((\textit{Spaces}/X)_\etale, \mathcal{O})$,
\item the rule $\mathcal{F} \mapsto \mathcal{F}^a$ defines
an equivalence between quasi-coherent $\mathcal{O}_X$-modules
and quasi-coherent modules on
$((\textit{Spaces}/X)_{fppf}, \mathcal{O})$,
\item we have $\epsilon_{X, *}a_X^*\mathcal{F} = \pi_X^*\mathcal{F}$
and $a_{X, *}a_X^*\mathcal{F} = \mathcal{F}$,
\item we have $R^i\epsilon_{X, *}(a_X^*\mathcal{F}) = 0$
and $R^ia_{X, *}(a_X^*\mathcal{F}) = 0$ for $i > 0$.
\end{enumerate}
\end{lemma}

\begin{proof}
Part (1) is a consequence of fppf descent of quasi-coherent modules.
Namely, suppose that $\{f_i : U_i \to U\}$ is an fpqc covering
in $(\textit{Spaces}/X)_\etale$. Denote $g : U \to X$ the structure
morphism. Suppose that
we have a family of sections $s_i \in \Gamma(U_i , f_i^*g^*\mathcal{F})$
such that $s_i|_{U_i \times_U U_j} = s_j|_{U_i \times_U U_j}$.
We have to find the correspond section $s \in \Gamma(U, g^*\mathcal{F})$.
We can reinterpret the $s_i$ as a family of maps
$\varphi_i : f_i^*\mathcal{O}_U = \mathcal{O}_{U_i} \to f_i^*g^*\mathcal{F}$
compatible with the canonical descent data associated to the
quasi-coherent sheaves $\mathcal{O}_U$ and $g^*\mathcal{F}$ on $U$.
Hence by Descent on Spaces, Proposition
\ref{spaces-descent-proposition-fpqc-descent-quasi-coherent}
we see that we may (uniquely) descend
these to a map $\mathcal{O}_U \to g^*\mathcal{F}$ which gives
us our section $s$.

\medskip\noindent
We will deduce (2) -- (7) from the corresponding statement for schemes.
Choose an \'etale covering $\{X_i \to X\}_{i \in I}$
where each $X_i$ is a scheme. Observe that $X_i \times_X X_j$
is a scheme too. This covering induces a covering of
the final object in each of the three sites
$(\textit{Spaces}/X)_{fppf}$, $(\textit{Spaces}/X)_\etale$, and $X_\etale$.
Hence we see that the category of sheaves on these sites
are equivalent to descent data for these coverings, see
Sites, Lemma \ref{sites-lemma-mapping-property-glue}.
Parts (2), (3) are local (because we have the glueing statement).
Being quasi-coherent is a local property, hence parts
(4), (5) are local. Clearly (6) and (7) are local.
It follows that it suffices to prove parts (2) -- (7)
of the lemma when $X$ is a scheme.

\medskip\noindent
Assume $X$ is a scheme. The embeddings
$(\Sch/X)_\etale \subset (\textit{Spaces}/X)_\etale$ and
$(\Sch/X)_{fppf} \subset (\textit{Spaces}/X)_{fppf}$
determine equivalences of ringed topoi by
Lemma \ref{lemma-compare-cohomology-other-topologies}.
We conclude that (2) -- (7) follows from the case of schemes.
\'Etale Cohomology, Lemma
\ref{etale-cohomology-lemma-review-quasi-coherent}.
To transport the property of being quasi-coherent via
this equivalence use that being quasi-coherent is an
intrinsic property of modules as explained in
Modules on Sites, Section \ref{sites-modules-section-local}.
Some minor details omitted.
\end{proof}

\section{Comparing big and small topoi}
\label{section-compare}

\begin{lemma}
\label{lemma-compare-cohomology-other-topologies}
Let $S$ be a scheme. Let $\tau \in \{\etale, fppf, ph\}$ (add more here).
The inclusion functor
$$
(\Sch/S)_\tau \longrightarrow (\textit{Spaces}/S)_\tau
$$
is a special cocontinuous functor
(Sites, Definition \ref{sites-definition-special-cocontinuous-functor})
and hence identifies topoi.
\end{lemma}

\begin{proof}
The conditions of Sites, Lemma \ref{sites-lemma-equivalence}
are immediately verified as our functor is fully faithful
and as every algebraic space has an \'etale covering by schemes.
\end{proof}

\begin{lemma}
\label{lemma-vanishing-adequate}
Let $S$ be a scheme. Let $X$ be an algebraic space over $S$.
Let $\mathcal{F}_1 \to \mathcal{F}_2 \to \mathcal{F}_3$
be a complex of quasi-coherent $\mathcal{O}_X$-modules.
Set
$$
\mathcal{H}_\etale =
\Ker(\pi_X^*\mathcal{F}_2 \to \pi_X^*\mathcal{F}_3)/
\Im(\pi_X^*\mathcal{F}_1 \to \pi_X^*\mathcal{F}_2)
$$
on $(\textit{Spaces}/X)_\etale$ and set
$$
\mathcal{H}_{fppf} =
\Ker(a_X^*\mathcal{F}_2 \to a_X^*\mathcal{F}_3)/
\Im(a_X^*\mathcal{F}_1 \to a_X^*\mathcal{F}_2)
$$
on $(\textit{Spaces}/X)_{fppf}$.
Then $\mathcal{H}_\etale = \epsilon_{X, *}\mathcal{H}_{fppf}$
and
$$
H^p_\etale(U, \mathcal{H}_\etale) = H^p_{fppf}(U, \mathcal{H}_{fppf}) = 0
$$
for $p > 0$ and any affine object $U$ of $(\textit{Spaces}/X)_\etale$.
\end{lemma}

\begin{proof}
For any object $f : U \to X$ of $(\textit{Spaces}/X)_\etale$
consider the restriction
$\mathcal{H}_\etale|_{U_\etale}$ of $\mathcal{H}_\etale$ to $U_\etale$ via
the functor $i_f^* = i_f^{-1}$ discussed in Section \ref{section-compare}.
The sheaf $\mathcal{H}_\etale|_{U_\etale}$
is equal to the homology of complex $f^*\mathcal{F}_\bullet$ in degree $1$.
This is true because $i_f \circ \pi_X = f$ as morphisms of ringed sites
$U_\etale \to X_\etale$. In particular we see that
$\mathcal{H}_\etale|_{U_\etale}$ is a quasi-coherent $\mathcal{O}_U$-module.
Next, let $g : V \to U$ be a flat
morphism in $(\textit{Spaces}/X)_\etale$. Since
$$
i_{f \circ g}^* \circ \pi_X^* = (f \circ g)^* = g^* \circ f^*
$$
as morphisms of sites $V_\etale \to X_\etale$ and since $g$ is flat
hence $g^*$ is exact, we obtain
$$
\mathcal{H}_\etale|_{V_\etale} =
g^*\left(\mathcal{H}_\etale|_{U_\etale}\right)
$$
With these preparations we are ready to prove the lemma.

\medskip\noindent
Let $\mathcal{U} = \{g_i : U_i \to U\}_{i \in I}$ be an fppf covering
with $f : U \to X$ as above. The sheaf property holds for
$\mathcal{H}_\etale$ and the covering $\mathcal{U}$
by (1) of Lemma \ref{lemma-review-quasi-coherent}
applied to $\mathcal{H}_\etale|_{U_\etale}$ and the above.
Therefore we see that $\mathcal{H}_\etale$ is already an fppf
sheaf and this means that $\mathcal{H}_{fppf}$ is
equal to $\mathcal{H}_\etale$
as a presheaf. In particular
$\mathcal{H}_\etale = \epsilon_{X, *}\mathcal{H}_{fppf}$.

\medskip\noindent
Finally, to prove the vanishing, we use
Cohomology on Sites, Lemma \ref{sites-cohomology-lemma-cech-vanish-collection}.
We let $\mathcal{B}$ be the affine objects of
$(\textit{Spaces}/X)_{fppf}$ and we let
$\text{Cov}$ be the set of finite fppf coverings
$\mathcal{U} = \{U_i \to U\}_{i = 1, \ldots, n}$ with $U$, $U_i$ affine.
We have
$$
{\check H}^p(\mathcal{U}, \mathcal{H}_\etale) =
{\check H}^p(\mathcal{U}, \left(\mathcal{H}_\etale|_{U_\etale}\right)^a)
$$
because the values of $\mathcal{H}_\etale$ on the
affine schemes $U_{i_0} \times_U \ldots \times_U U_{i_p}$
flat over $U$ agree with the values of the pullback
of the quasi-coherent module $\mathcal{H}_\etale|_{U_\etale}$ by
the first paragraph. Hence we obtain vanishing by
Descent, Lemma \ref{descent-lemma-standard-covering-Cech-quasi-coherent}.
This finishes the proof.
\end{proof}

\end{document}
